% Last updated: 2025/04/22

% Main specification
Upon graduating from residency and fellowship, physicians migrate out of necessity. There will be those who begin practice in a different market than the one they trained in, who I call ``movers", and those who stay in the market they trained in, who I call ``stayers". The core of my identification strategy is to compare physician movers to stayers while controlling for influences from their current market. This utilizes the fact that movers who train in a different market are exposed to a different intensity level in medical school compared to stayers. The primary specification is 
    \begin{align} \label{eq:main}
        intensity_{it} 
        &= \beta (mover_i \times \textit{med school HRR intensity}_{it_g}) \nonumber\\
        &\quad + \psi \textit{med school HRR intensity}_{it_g} + \theta X_{i} + \eta M_{it} \nonumber\\
        &\quad + \gamma_i + \tau_t + \epsilon_{it},
    \end{align}
where $intensity_{it}$ is physician $i$'s treatment intensity in year $t$, $mover_i$ is an indicator that physician $i$'s current practice HRR is different from their medical school HRR, and $\textit{med school HRR intensity}_{it_g}$ is the average treatment intensity in physician $i$'s medical school HRR in year of graduation $t_g$. The vector $X_i$ contains physician characteristics, and the vector $M_{it}$ contains variables that reflect the physician's Medicare patient population, as described above. All models include current HRR fixed effects $\gamma_i$ and year fixed effects $\tau_t$. Standard errors are clustered at the current HRR level due to spacial correlation. 

% Semi-parametric specification
I estimate equation \ref{eq:main} using OLS, which assumes that the effect of medical school intensity on a physician's current intensity is linear. However, one can easily imagine that more extreme medical school intensities will have a greater impact. For this reason, I also semi-parametrically estimate equations of the following form: 
    \begin{align} \label{eq:spline}
        intensity_{it} 
        &= f(\textit{med school HRR intensity}_{it_g}) \times mover_i \nonumber\\
        &\quad + \theta X_{i} + \eta M_{it} + \gamma_i + \tau_t + \epsilon_{it},
    \end{align}
where $f(\cdot)$ is estimated using a cubic spline. 

These two specifications capture the effect of the level of medical school intensity. However, they fail to acknowledge that some movers now practice in more intense HRRs compared to their medical school, while others practice in less intense HRRs compared to their medical school. To account for this fact, I additionally estimate equations using the difference between current HRR intensity and medical school intensity. 
    \begin{align} \label{eq:delta}
        intensity_{it} 
        &= \beta (mover_i \times \textit{change in intensity}_{it}) \nonumber\\
        &\quad + \psi \textit{med school HRR intensity}_{it_g} + \theta X_{i} + \eta M_{it} \nonumber\\
        &\quad + \gamma_i + \tau_t + \epsilon_{it},
    \end{align}
where $\textit{change in intensity}_{it}$ is defined as the difference between the leave-one-out mean intensity of physician $i$'s current HRR in year $t$, excluding physician $i$, and the average treatment intensity in physician $i$'s medical school HRR in their year of graduation. 

% Endogeneity concerns
There are two potential sources for endogeneity: reverse causality and omitted variable bias. It is likely that average HRR intensities are changing over time due to physician movement. Since my medical school intensity measures are calculated at the time the physician graduates from medical school, reverse causality is not an issue. Omitted variable bias is a greater concern here. It is possible that medical schools are investing in new technology, and this improvement in technology is what is actually responsible for both changes in market-level treatment intensity as well as the intensity of graduating physicians. It may be possible to account for this using AHA survey data and requires further analysis. 
